%% -*- coding: utf-8 -*-

\chapter{词性}

\section{名词和冠词}

\subsection{名词类别以及复合名词}

\begin{center}
  \begin{talltblr}[
    label = none, % Removes the "Table X:" text
    ]{
    vlines, hlines,
    % 加入支持列表环境的测量参数 (记得在导言区加载 \UseTblrLibrary{varwidth})
    % measure=vbox, stretch=-1,
    rows={halign=c, valign=m},
    columns={halign=c, valign=m},
    column{1}={0.25\linewidth, halign=c, bg=azure9},
    column{2}={0.6\linewidth, halign=l},
    % column{3}={0.35\linewidth, halign=l},
    % column{4}={0.35\linewidth, halign=l},
    row{1}={font=\large\bfseries, bg=azure9, halign=c}
    }
    形式                     & 举例                                                                                                         \\
    两个或多个词连写         & handbag(\(n\). + \(n\).), greenhouse(\(adj\). + \(n\).), output(\(adv\). + \(v\).), underground(\(prep\). + \(n\).), handwriting(\(n\). + \(v\).ing) \\
    两个或多个词用连字符连接 & good-looking(\(adj\).+\(v\).ing), fire-raiser(\(n\).+\(n\).), life-long(\(n\).+\(adj\).), summing-up(\(v\).ing+\(adv\).), mother-in-law(\(n\).+\(prep\).+\(n\).)\\
    两个或多个词分开写       & river bank(\(n\).+\(n\).), driving license(\(v\).ing+\(n\).)\\
  \end{talltblr}
\end{center}

\markword[\uuline]{\textbf{Britain}}{主语 1}
\markword{had produced}{谓语}
\markword[\uwave]{textiles}{宾语}
\markword[\mbox]{(like wool, linen and cotton),}{后置定语 (修饰 textiles)}
\markword[\mbox]{[for hundreds of years],}{时间状语}

\vspace{0.5cm} % 手动换行并增加行距

\markword[\mbox]{but}{连词}
\markword[\mbox]{[prior to the Industrial Revolution],}{时间状语}
\markword[\uuline]{\textbf{the British textile business}}{主语 2}
\markword{was}{系动词}

\vspace{0.5cm}

\markword[\uuwave]{a true ``cottage industry'',}{表语}
\markword[\mbox]{[with the work performed in small workshops or even homes by}{伴随状语}

\vspace{0.5cm}

individual spinners, weavers and dyers].

\subsection{可数名词变复数（规则变化）}

名词除了根据意义划分，还可以根据其可数性分为“可数名词”和“不可数名词”。其中可数名词在变复数时，一般情况下在词尾加-s或-es，如句中的tardigrades、temperatures和degrees都是可数名词的复数形式，直接在词尾加-s。
\begin{center}
  \begin{talltblr}[
    label = none, % Removes the "Table X:" text
    ]{
    vlines, hlines,
    % 加入支持列表环境的测量参数 (记得在导言区加载 \UseTblrLibrary{varwidth})
    % measure=vbox, stretch=-1,
    rows={halign=c, valign=m},
    columns={halign=c, valign=m},
    column{1}={0.1\linewidth, halign=c, bg=azure9},
    column{2}={0.35\linewidth, halign=l},
    column{3}={0.4\linewidth, halign=l},
    % column{4}={0.35\linewidth, halign=l},
    row{1}={font=\large\bfseries, bg=azure9, halign=c}
    }
    类别&构成方法&举例\\
    \SetCell[r=3]{c}结尾加-s&词尾直接加-s&restaurant \(\rightarrow\) restaurants; professor \(\rightarrow\) professors\\
        &以“元音+y”结尾的名词，加-s&attorney \(\rightarrow\) attorneys; alloy \(\rightarrow\) alloys\\
        &以“元音+o”结尾的名词，加-s&studio \(\rightarrow\) studios; kangaroo \(\rightarrow\) kangaroos\\
    \SetCell[r=4]{c}结尾加-es&以 [s]、[z]、[ʃ]、[ʒ]、[ʧ]、[ʤ]等音节结尾的名词，加-es&buzz \(\rightarrow\) buzzes; bunch \(\rightarrow\) bunches; brush \(\rightarrow\) brushes; watch \(\rightarrow\) watches\\
        &以“辅音+o”结尾的名词，加 -es&volcano \(\rightarrow\) volcanoes; potato \(\rightarrow\) potatoes\\
        &词尾为-f或-fe，变f、fe 为v，加-es&shelf \(\rightarrow\) shelves; thief \(\rightarrow\) thieves\\
        &以“辅音+y”结尾的名词，变y为i加-es&strategy \(\rightarrow\) strategies; agency \(\rightarrow\) agencies\\
  \end{talltblr}
\end{center}


\markword[\mbox]{There's been}{there be 结构}
\markword[\uuline]{\textbf{some interesting research}}{主语}
\markword[\mbox]{{\large\textcolor{red}{\(\lgroup\)}}}{\textcolor{red}{定从(research)开始}}
\markword[\mbox]{which}{引导词/主语}

\vspace{0.5cm}

\markword{has found}{谓语}
\markword[\mbox]{\textcolor{blue}{\(\lgroup\)}}{\textcolor{blue}{宾从开始}}
\markword[\mbox]{that}{引导词}
\markword[\uuline]{tardigrades}{主语 1}
\markword{are capable of}{谓语 1}
\markword[\uwave]{surviving radiation and very high pressure,}{宾语 1}

\vspace{0.5cm}

\markword[\mbox]{and}{连词}
\markword[\uuline]{they}{主语 2}
\markword{are also able to withstand}{谓语 2}
\markword[\uwave]{temperatures}{宾语 2}
\markword[\mbox]{(as cold as -200 degrees centigrade),}{后置定语（temperatures）}

\vspace{0.5cm}

\markword[\mbox]{or}{连词}
\markword[\uwave]{highs}{宾语 3}
\markword[\mbox]{(of more than 148 degrees centigrade),}{后置定语（highs）}
\markword[\mbox]{(which is incredibly hot).}{定语从句}
\markword[\mbox]{\textcolor{blue}{\(\rgroup\)}}{\textcolor{blue}{宾从结束}}
\markword[\mbox]{\large{\textcolor{red}{\(\rgroup\)}}}{\textcolor{red}{定从结束}}

\subsection{可数名词变复数（不规则变化）}

可数名词变复数并非都是规则变化，还有一些特殊变化，比如句中的phenomena就是复数形式，其原形为phenomenon。
\begin{center}
  \begin{talltblr}[
    label = none, % Removes the "Table X:" text
    ]{
    vlines, hlines,
    % 加入支持列表环境的测量参数 (记得在导言区加载 \UseTblrLibrary{varwidth})
    % measure=vbox, stretch=-1,
    rows={halign=c, valign=m},
    columns={halign=c, valign=m},
    column{1}={0.25\linewidth, halign=c, bg=azure9},
    column{2}={0.55\linewidth, halign=l},
    % column{3}={0.35\linewidth, halign=l},
    % column{4}={0.35\linewidth, halign=l},
    row{1}={font=\large\bfseries, bg=azure9, halign=c}
    }
    构成方法&举例\\
    改变中间元音&goose \(\rightarrow\) geese; policeman \(\rightarrow\) policemen; analysis \(\rightarrow\) analyses; tooth \(\rightarrow\) teeth\\
    单复数同形&sheep \(\rightarrow\) sheep; salmon \(\rightarrow\) salmon\\
    结尾加-en或-ren&ox \(\rightarrow\) oxen; child \(\rightarrow\) children\\
    其他不规则形式&phenomenon \(\rightarrow\) phenomena; criterion \(\rightarrow\) criteria\\
  \end{talltblr}
\end{center}


\markword[\mbox]{[In the 1960s],}{时间状语}
\markword[\uuline]{the astronomer Gerald Hawkins}{主语}
\markword{suggested}{谓语}
\markword[\mbox]{\large\textcolor{blue}{\(\lgroup\)}}{宾从开始}
\markword[\mbox]{that}{引导词}

\vspace{0.5cm}

\markword[\uuline]{the cluster of megalithic stones}{主语}
\markword[\mbox]{operated}{谓语}
\markword{[as a form of calendar],}{方式状语}

\vspace{0.5cm}

\markword[\mbox]{[with different points corresponding to astronomical phenomenal]}{伴随状语}

\vspace{0.5cm}

\markword[\mbox]{<such as solstices, equinoxes and eclipses>}{同位语（修饰 phenomenal）}
\markword[\mbox]{(occurring at different times of the year).}{定语（修饰同位语）}
\markword[\mbox]{\large\textcolor{blue}{\(\rgroup\)}}{宾从结束}

\subsection{不可数名词}

不可数名词是指无法数清具体数量的名词，常见的不可数名词分为以下四类。
\begin{center}
  \begin{talltblr}[
    label = none, % Removes the "Table X:" text
    ]{
    vlines, hlines,
    % 加入支持列表环境的测量参数 (记得在导言区加载 \UseTblrLibrary{varwidth})
    % measure=vbox, stretch=-1,
    rows={halign=c, valign=m},
    columns={halign=c, valign=m},
    column{1}={0.15\linewidth, halign=c, bg=azure9},
    column{2}={0.35\linewidth, halign=l},
    column{3}={0.35\linewidth, halign=l},
    % column{4}={0.35\linewidth, halign=l},
    row{1}={font=\large\bfseries, bg=azure9, halign=c}
    }
    不可数类别&举例&不可数情况说明\\
    分割不变&wine, butter, honey&切开这类物质后，性质没有变化\\
    难以数清&dust, flour, pepper, hair&数量很多，很难数清楚\\
    总称概念&data, furniture, clothing&表示总称的概念，不是一个具体的事物\\
    抽象名词&health, happiness, patience, ambition, knowledge, courage&抽象的思想、观点等，无法计数\\
  \end{talltblr}
\end{center}
注意:还有一些不可数名词看上去像是可数名词的复数形式，但实际上却是不可数名词，如本句中的economics，类似的词还有news 以及 physics。

\markword[\uuline]{The book}{主语}
\markword{combines}{谓语}
\markword[\uwave]{geology, history, economics, and a lot of data}{宾语}

\vspace{0.5cm}

\markword[\mbox]{\Large\textcolor{orange}{\( [ \)}}{\textcolor{orange}{目的壮语开始}}
\markword[\mbox]{to explain}{}
\markword[\mbox]{why}{引导词}
\markword[\uuline]{business clusters}{主语}
\markword{developed}{谓语}
\markword[\mbox]{[where they did]}{地点状语从句}

\vspace{0.5cm}

\markword[\mbox]{how}{引导词}
\markword[\uuline]{the early decisions of workers and firms}{主语}
\markword{shaped}{谓语}
\markword[\uwave]{the skyline}{宾语}
\markword[\mbox]{(we see today).}{定语从句}
\markword[\mbox]{\Large\textcolor{orange}{\( ] \)}}{\textcolor{orange}{目的壮语结束}}

\subsection{名词修饰语}

想要描述不同类型名词的数量时需要注意，有些修饰语只能修饰可数名词，有些只能修饰不可数名词，有些二者均可修饰。如a picce of在本句中就用来修饰不可数名词art，表示“一件艺术品”。常用的名词修饰语还有:
\begin{center}
  \begin{talltblr}[
    label = none, % Removes the "Table X:" text
    ]{
    vlines, hlines,
    % 加入支持列表环境的测量参数 (记得在导言区加载 \UseTblrLibrary{varwidth})
    % measure=vbox, stretch=-1,
    rows={halign=c, valign=m},
    columns={halign=c, valign=m},
    column{1}={0.2\linewidth, halign=c, bg=azure9},
    column{2}={0.65\linewidth, halign=l},
    % column{3}={0.35\linewidth, halign=l},
    % column{4}={0.35\linewidth, halign=l},
    row{1}={font=\large\bfseries, bg=azure9, halign=c}
    }
    类别&举例\\
    只修饰可数名词&few (很少，几乎没有), several (几个), many (许多), a number of (大量的), scant (为数不多的，不足的), numerous (众多的，许多的), a handful of (少数，几个), a few (几个)\\
    只修饰不可数名词&little (很少，几乎没有), a little (一些，少量), a bit of (有点儿), much (许多), a large amount of (大量的), a great deal of (很多), a great quantity of (大量的), an abundance of (丰富的，充裕的)\\
    可数名词、不可数名词均可修饰&some (一些), plenty of (许多), most (大多数的), a wide range of (广泛的), substantial (大量的), a lot of/lots of (许多), enough (足够的), various (各种各样的), ample (充足的), considerable (相当多的，相当大的)\\
  \end{talltblr}
\end{center}


\markword[\mbox]{[Likewise],}{状语}
\markword[\mbox]{\Large\textcolor{orange}{\( [ \)}}{\textcolor{orange}{条件壮语开始}}
\markword[\mbox]{if}{引导词}
\markword[\uuline]{the tree}{主语}
\markword{suffers from}{谓语}
\markword[\uwave]{black knot disease}{宾语}
\markword[\mbox]{{\Large\textcolor{orange}{\( ] \)}},}{\textcolor{orange}{条件壮语结束}}

\vspace{0.5cm}

\markword[\uuline]{its value for timber}{主语 1}
\markword{decreases,}{谓语}
\markword[\mbox]{but}{连词}
\markword[\mbox]{[to a woodworker]}{对象状语}
\markword[\mbox]{(interested in making bowls),}{后置定语（woodworker）}

\vspace{0.5cm}

\markword[\uuline]{it}{主语 2}
\markword{brings}{谓语 2}
\markword[\uwave]{an opportunity}{宾语}
\markword[\mbox]{(for a unique and beautiful piece of art).}{后置定语（opportunity）}

\subsection{即可数又不可数名词}

有些名词既可数又不可数，含义会随语境而变，比如:
\begin{center}
  \begin{talltblr}[
    label = none, % Removes the "Table X:" text
    ]{
    vlines, hlines,
    % 加入支持列表环境的测量参数 (记得在导言区加载 \UseTblrLibrary{varwidth})
    % measure=vbox, stretch=-1,
    rows={halign=c, valign=m},
    columns={halign=c, valign=m},
    column{1}={0.15\linewidth, halign=c, bg=azure9},
    column{2}={0.35\linewidth, halign=l},
    column{3}={0.35\linewidth, halign=l},
    % column{4}={0.35\linewidth, halign=l},
    row{1}={font=\large\bfseries, bg=azure9, halign=c}
    }
    名词&不可数&可数\\
    paper&{含义:纸\\I need some paper to write on. 我需要一些纸来写字。}&{含义:论文;试卷\\These three papers on artificial intelligence are really insightful. 这三篇关于人工智能的论文很有见地。}\\
    room&{含义:空间\\There is no room for you to put another chair here. 这里没有空间让你再放一把椅子了。}&{含义:房间\\There are five rooms in this house.这所房子有五个房间。}\\
    light&{含义:光;光线\\The light in the room is too dim. 房间里的光线太暗了。}&{含义:灯;光源\\There are several lights in the street. 街上有几盏灯。}\\
  \end{talltblr}
\end{center}


\markword[\uuline]{The burial chamber of the tomb,}{主语}
\markword[\mbox]{{\Large\textcolor{red}{\( \lgroup \)}}}{\textcolor{red}{定从1(tomb)开始}}
\markword[\mbox]{where}{引导词}
\markword[\uuline]{the king's body}{主语}
\markword{was laid}{谓语}

\vspace{0.5cm}

\markword[\mbox]{[to rest],}{目的状语}
\markword[\mbox]{{\Large\textcolor{red}{\( \rgroup \)}}}{\textcolor{red}{定从1结束}}
\markword{was dug}{谓语}
\markword[\mbox]{[beneath the base of the pyramid],}{地点状语}

\vspace{0.5cm}

\markword[\mbox]{(surrounded by a vast maze of long tunnels)}{后置定语（chamber）}
\markword[\mbox]{{\Large\textcolor{red}{\( \lgroup \)}}}{\textcolor{red}{定从2(tunnels)开始}}
\markword[\mbox]{that}{引导词/主语}


\vspace{0.5cm}

\markword{had}{谓语}
\markword[\uwave]{rooms}{宾语}
\markword[\mbox]{[off them]}{状语}
\markword[\mbox]{[to discourage robbers].}{目的状语}
\markword[\mbox]{{\Large\textcolor{red}{\( \rgroup \)}}}{\textcolor{red}{定从2结束}}

\subsection{名词所有格}

在表达所属关系时，可以借助名词所有格。名词所有格主要分为's所有格和of所有格，有时两者会一同使用，即“双重所有格”。
\begin{center}
  \begin{talltblr}[
    label = none, % Removes the "Table X:" text
    ]{
    vlines, hlines,
    % 加入支持列表环境的测量参数 (记得在导言区加载 \UseTblrLibrary{varwidth})
    % measure=vbox, stretch=-1,
    rows={halign=c, valign=m},
    columns={halign=c, valign=m},
    column{1}={0.1\linewidth, halign=c, bg=azure9},
    column{2}={0.1\linewidth, halign=c, bg=azure9},
    column{3}={0.2\linewidth, halign=l},
    column{4}={0.45\linewidth, halign=l},
    row{1}={font=\large\bfseries, bg=azure9, halign=c}
    }
    类型&词尾&构成&举例\\
    \SetCell[r=2]{c}'s所有格&一般性词尾&词尾+'s&the scientist's research findings (科学家的研究成果)\\
        &-s, -es&词尾+'&two hours’ drive from the city center (离市中心两小时车程的地方)\\
    \SetCell[c=2]{c}of所有格&&名词+of+名词&the complexity of the economic system (经济体系的复杂性)\\
    \SetCell[r=2,c=2]{c}双重所有格&&of+名词-'s&a painting of Picasso’s (毕加索的一幅画作)\\
        &&of+名词性物主代词&a novel of his (他的一本小说)\\
  \end{talltblr}
\end{center}


\markword[\mbox]{[By the mid-1920s],}{时间状语}
\markword[\uuline]{O'keeffe}{主语}
\markword{was recognized}{谓语}

\vspace{0.5cm}

\markword[\mbox]{<as one of America's most important and successful artists>,}{主语补足语}

\vspace{0.5cm}

\markword[\mbox]{(widely known for the architectural pictures)}{后置定语（主语）}
\markword[\mbox]{{\Large\textcolor{red}{\( \lgroup \)}}}{\textcolor{red}{定从(pictures)开始}}

\vspace{0.5cm}

\markword[\mbox]{that}{引导词/主语}
\markword[\mbox]{[dramatically]}{状语}
\markword{depict}{谓语}
\markword[\uwave]{the soaring skyscrapers of New York.}{宾语}
\markword[\mbox]{{\Large\textcolor{red}{\( \rgroup \)}}}{\textcolor{red}{定从结束}}

\subsection{不定冠词a和an}

不定冠词有a和 an,a用在辅音音素前面（注意不是辅音字母），如a challenge“一项挑战”；an用在元音音素前面（注意不是元音字母），如本句中的an era“一个时代”。
\begin{center}
  \begin{talltblr}[
    label = none, % Removes the "Table X:" text
    ]{
    vlines, hlines,
    % 加入支持列表环境的测量参数 (记得在导言区加载 \UseTblrLibrary{varwidth})
    % measure=vbox, stretch=-1,
    rows={halign=c, valign=m},
    columns={halign=c, valign=m},
    column{1}={0.1\linewidth, halign=c, bg=azure9},
    column{2}={0.23\linewidth, halign=l},
    column{3}={0.55\linewidth, halign=l},
    % column{4}={0.35\linewidth, halign=l},
    row{1}={font=\large\bfseries, bg=azure9, halign=c}
    }
    用法&含义&举例\\
    \SetCell[r=5]{c}一般用法&泛指某一类人、事或物&A dog is a loyal animal. 狗是一种忠诚的动物。\\
        &首次提到某人或某物时&I saw a man walking down the street. 我看到一个男人在街上走。\\
        &表示数量“一”&Can I have a cup of coffee, please? 请给我一杯咖啡好吗?\\
        &表示“每一”&The car is running at a speed of 80 kilometers an hour. 汽车以每小时 80千米的速度行驶。\\
        &表示“又一，再一”&Can you give me a second chance? 你能再给我一次机会吗?\\
    特殊用法&用于抽象名词前&{surprise (惊讶) \(\rightarrow\) a surprise (一件令人惊讶的事)\\succes(成功) \(\rightarrow\) a success (一个成功的人;一件成功的事)}\\
    固定搭配&\SetCell[c=2]{l}{all of a sudden (突然地，出乎意料地), in a hurry (立即，匆忙),\\as a result of (由于), at a loss (困惑不解，茫然不知所措),\\play a key/vital/crucial role in (在……中起关键作用)}\\
  \end{talltblr}
\end{center}


\markword[\uuline]{\textbf{We}}{主语 1}
\markword[\mbox]{[currently]}{时间状语}
\markword{live}{谓语}
\markword[\mbox]{[in an era]}{时间状语}
\markword[\mbox]{{\Large\textcolor{red}{\( \lgroup \)}}}{\textcolor{red}{定从(an era)开始}}
\markword[\mbox]{in which}{引导词}
\markword[\uuline]{\textbf{technology}}{主语}
\markword{enables}{谓语}

\vspace{0.5cm}

\markword[\uwave]{information}{宾语}
\markword[\mbox]{<to reach large audiences>}{主语补足语}
\markword[\mbox]{(distributed across the globe),}{后置定语（audiences）}
\markword[\mbox]{and thus}{连词}

\vspace{0.5cm}

\markword[\uuline]{\textbf{the potential}}{主语 2}
\markword[\mbox]{(for immediate and widespread effects from misinformation)}{后置定语（potential）}

\vspace{0.5cm}

\markword[\mbox]{[now]}{时间状语}
\markword{looms}{系动词}
\markword[\uuwave]{larger}{表语}
\markword[\mbox]{[than in the past].}{比较状语}
\markword[\mbox]{{\Large\textcolor{red}{\( \rgroup \)}}}{\textcolor{red}{定从结束}}

\subsection{定冠词the}

定冠词the用法广泛，如用于序数词前（本句中的 the first 和 the third），用于专有名词前（本句中的 the Caribbean 和the Mediterranean）。现将定冠词的用法汇总如下:
\begin{center}
  \begin{talltblr}[
    label = none, % Removes the "Table X:" text
    ]{
    vlines, hlines,
    % 加入支持列表环境的测量参数 (记得在导言区加载 \UseTblrLibrary{varwidth})
    % measure=vbox, stretch=-1,
    rows={halign=c, valign=m},
    columns={halign=c, valign=m},
    column{1}={0.1\linewidth, halign=c, bg=azure9},
    column{2}={0.1\linewidth, halign=l},
    column{3}={0.3\linewidth, halign=l},
    column{4}={0.35\linewidth, halign=l},
    row{1}={font=\large\bfseries, bg=azure9, halign=c}
    }
    用法&指向&含义&举例\\
    \SetCell[r=11]{c}一般用法&\SetCell[r=2]{c}表特指&特指上文提到的人或事物&第一次提到用 a/an，以后再次提到用 the\\
        &&特指谈话双方都明确的人或物&Close the window, please. 请关上窗户。\\
        &\SetCell[r=5]{c}专用名词前&党派、民族或国家名称等专有名词&the United States 美国\\
        &&海洋、江河、湖泊、山脉、海峡、海湾等地理名词&the Yellow River 黄河\\
        &&世界上、宇宙中独一无二的事物&the earth 地球\\
        &&表示演奏乐器时，用于乐器前&play the erhu 拉二胡\\
        &&著名的历史建筑&the Great Wall 长城\\
        &\SetCell[r=2]{c}时间和方位&用于世纪或年代前&in the 1980s在20世纪 80年代\\
        &&用于方位词前面&the west 西方\\
        &\SetCell[c=2]{l}用在序数词和形容词最高级前&&the most beautiful 最漂亮的\\
        &\SetCell[c=2]{l}用于姓氏复数形式前，表示“全家人”或“夫妇俩”&&the Greens格林一家\\
    \SetCell[r=2]{c}特殊用法&\SetCell[r=2]{c}the+形容词&表示一类人（表示复数）&the poor 穷人\\
        &&表示抽象名词 (表示单数)&the right 正确的事/做法\\
    固定搭配&\SetCell[c=3]{l}{with the passage of time 随着时间的流逝\qquad{}on the grounds of 由于\\at the high expense of 以……为代价}\\
  \end{talltblr}
\end{center}


\markword[\mbox]{\Large\textcolor{orange}{\( [ \)}}{\textcolor{orange}{时壮开始}}
\markword[\mbox]{Yet from the first to the third millennium BCE, thousands of years}{}

\vspace{0.5cm}

\markword[\mbox]{\large\textcolor{olive}{\( [ \)}}{\textcolor{olive}{壮从开始}}
\markword[\mbox]{before}{引导词}
\markword[\uuline]{\textbf{these swashbucklers}}{主语}
\markword{began}{谓语}
\markword[\uwave]{spreading fear}{宾语}
\markword[\mbox]{[across the Caribbean]}{地点状语}

\vspace{0.5cm}

\markword[\mbox]{{\large\textcolor{olive}{\( ] \)}}}{\textcolor{olive}{壮从结束}}
\markword[\mbox]{{\Large\textcolor{orange}{\( ] \)}},}{\textcolor{orange}{时壮结束}}
\markword[\uuline]{\textbf{pirates}}{主语}
\markword{prowled}{谓语}
\markword[\uwave]{the Mediterranean,}{宾语}

\vspace{0.5cm}

\markword[\mbox]{[raiding merchant ships and threatening vital trade routes].}{伴随状语}

\subsection{零冠词}

有些情况下，名词前不需要添加冠词，即“零冠词”，如本句中的 at least(至少)。
\begin{center}
  \begin{talltblr}[
    label = none, % Removes the "Table X:" text
    ]{
    vlines, hlines,
    % 加入支持列表环境的测量参数 (记得在导言区加载 \UseTblrLibrary{varwidth})
    % measure=vbox, stretch=-1,
    rows={halign=c, valign=m},
    columns={halign=c, valign=m},
    column{1}={0.1\linewidth, halign=c, bg=azure9},
    column{2}={0.3\linewidth, halign=l},
    column{3}={0.45\linewidth, halign=l},
    % column{4}={0.35\linewidth, halign=l},
    row{1}={font=\large\bfseries, bg=azure9, halign=c}
    }
    用法&含义&举例\\
    \SetCell[r=5]{c}一般用法&复数名词、不可数名词前&I need help with my group assignment. 我的小组作业需要帮助。\\
        &大陆、单一的国家名称、州县、城、镇等地名前&in Australia 在澳大利亚\\
        &人名、称呼和头衔前&Sorry,sir. 抱歉，先生。\\
        &季节、月份、日期、星期几、节假日前&Spring is the best season of the year. 春天是一年中最好的季节。\\
        &球类和棋类运动、学科、语言、三餐和疾病前&I have lunch at 2 p.m. 我下午两点吃午饭。\\
    固定搭配&\SetCell[c=2]{l}{at first (起初)\qquad{}in time (及时)\qquad{}in danger (在危险中)\\lose weight (减肥)\qquad{}on time (按时)\qquad{}by car (乘坐小汽车)}\\
  \end{talltblr}
\end{center}


\markword[\uuline]{\textbf{A review}}{主语}
\markword[\mbox]{(of almost 1,800 cases)}{后置定语（review）}
\markword[\mbox]{(of entanglement in fishing nets}{}

\vspace{0.5cm}

\markword[\mbox]{and of plastic consumption among marine mammals)}{后置定语（cases）}
\markword[\mbox]{[in US waters]}{地点状语}
\markword[\mbox]{[from 2009 to 2020]}{时间状语}

\vspace{0.5cm}

\markword{found}{谓语}
\markword[\mbox]{\large\textcolor{blue}{\(\lgroup\)}}{\textcolor{blue}{宾从开始}}
\markword[\mbox]{that}{引导词}
\markword[\uuline]{\textbf{at least 700 cases}}{主语}
\markword{involved}{谓语}
\markword[\uwave]{manatees}{宾语}
\markword[\mbox]{\large\textcolor{blue}{\(\rgroup\)}}{\textcolor{blue}{宾从结束}}


%%% Local Variables:
%%% TeX-engine: xetex
%%% TeX-master: "../IELTS_300_Sentences"
%%% End:

% LocalWords:  vlines hlines halign valign xetex
